\scanpage{32}
conjugaison de scindages. Donc il suffit maintenant de prouver le

\uline{Théorème 3} Soit $X$ un schéma de type fini sur un corps $K$, extension de type fini de $\mathbb Q$, on suppose que $X$ est \deleted{\unsure{connexe} lisse} \inserted{géom. 0-connexe,} et se plonge dans une extension d'une VA par un tore (p.ex. $X$ est une variété élémentaire d'Artin, à fibres types \inserted{fibres} distinctes de la droite projective et de la droite affine).

\margin{il suffirait $K$ de type fini sur un corps premier, pas néc. de car. 0 -- mais sous cette forme ce n'est pas très utilisable en car $p>0$, cf commentaires plus bas...}

Considérons \deleted{l'application} \inserted{une} clôture alg. $\overline K/K$, et \deleted{\unsure{la}} l'extension extérieure correspondante \inserted{$E_{X/K}$} de $E_{\overline K/K}=\operatorname{Gal}(\overline K/K)$ par $\pi_1(\overline X)$ $(\overline X=X\otimes_K\overline K)$, et l'extension déduite $\widetilde E_{X/K}$ de $E_{\overline K/K}$ par $\pi_1(\overline X)_{\mathrm{ab}}$. Considérons les applications canoniques
\diagram{\begin{tikzcd}[ampersand replacement=\&,column sep=large,row sep=large]
(*)\quad X(K) \arrow[r] \& {\begin{gathered}\text{classes de }\inserted{\pi(\overline X)\text{-}}\text{conjugaison de}\\\text{scindages de }E_{X/K}\text{ sur }E_{\overline K/K}\end{gathered}} \arrow[d,dotted] \\
(**)\quad X(K) \arrow[r] \& {\begin{gathered}\text{classes de conjugaison de scindages}\\\text{de l'ext. }\widetilde E_{X/K}\text{ de }E_{\overline K/K}\text{ par }H_1(Y).\end{gathered}}
\end{tikzcd}}
\margin{31}
